	
\section{Analoge Filter}

	\begin{symbolbox}
		$ f_c = \text{Grenzfrequenz (cutoff frequency)} $ \\
		$ \omega = \text{Kreisfrequenz}, 2\pi f $
	\end{symbolbox}

$\underline{H}$


\subsection{Passive Filter 1. Ordnung}

	6dB/Oktave bzw. 20dB/Dekade

	\subsection{RC-Filter}
	
	
	\begin{emphbox}
		\begin{align*}
			U_{2} &=U_{1}\cdot {\frac {\vert X_{C}\vert }{\sqrt {R^{2}+X_{C}^{2}}}}={\frac {U_1}{\sqrt {1+({\frac {f}{f}}_{c})^{2}}}} \\
		\end{align*}
	\end{emphbox}
	
	
		\begin{emphbox}
			Grenzfrequenz: $ f_c = \frac{1}{2 \pi RC} $ \\
			Pol: $ s = - \frac{1}{RC} $ \\
			$ H(s) = \frac{U_2}{U_1} $ \\

			\begin{tabular}{l|l}
				Tiefpass: & Hochpass : \\
				\subfile{circuits/rc_low_pass.tikz.tex} & 
				\subfile{circuits/rc_high_pass.tikz.tex} \\
				$ \frac{U_2}{U_1} = \frac{1}{\sqrt{1+(\omega RC)^2}} $ &
				$ \frac{U_2}{U_1} = \frac{R}{\sqrt{R^2+(\frac{1}{\omega C})^2}} $ \\
				$ H(s) = \frac{1}{1 + RCs} $ &
				$ H(s) = \frac{RCs}{1 + RCs} $
			\end{tabular}
		\end{emphbox}

	\subsection{RL-Filter}



		\begin{emphbox}
			
			
			Grenzfrequenz: $ f_c = \frac{R}{2 \pi L} $ \\
			Pol: $ s = - \frac{R}{L} $ \\
			Nullstelle: $ Origin $	\\
			$ H(s) = \frac{U_2}{U_1} $ \\
			
			\begin{tabular}{l|l}
				Tiefpass: & Hochpass: \\
				\subfile{circuits/rl_low_pass.tikz.tex} & 
				\subfile{circuits/rl_high_pass.tikz.tex} \\
				$ \frac{U_2}{U_1} = \frac{1}{\sqrt{1+(\frac{\omega L}{R})^2}} $ &
				$ \frac{U_2}{U_1} = \frac{1}{\sqrt{1+(\frac{R}{\omega L})^2}} $ \\
				$ H(s) = \frac{R}{R + Ls} $ &
				$ H(s) = \frac{Ls}{R + Ls} $
			\end{tabular}
		\end{emphbox}

\subsection{Digitale Filter}
\subsubsection{FIR-Filter}
Finite impulse filter, auch Transversalfilter.
Eigenschaften:
- nicht schwingungsfähig, daher stabil
- normalerweise nicht rekursiv
- meistens digital
- Polverteilung: n-fache Polstelle im Ursprung
- bei rekursiven FIR-Filtern fallen Nullstellen mit Polstellen zusammen -> stabil
- Koeffizientenzahl = Ordnung + 1

	\subsubsection{Impulsantwort}

	$g(t) = \begin{cases}
		0 & t < 0 \\
		\frac{K}{T}\cdot e^{-\frac{t}{T}} & t \ge 0
	\end{cases}$
	$g(t) = \rho(t)$
	
%g(t) Impulsantwort auf Einheitsimpuls \delta(t)

\begin{sectionbox}

\end{sectionbox}
